\documentclass[11pt,a4paper]{article}

\usepackage[utf8]{inputenc}
\usepackage[T1]{fontenc}
\usepackage[ngerman]{babel}
\usepackage{amsmath,amssymb,amsfonts}
\usepackage{physics}
\usepackage{mathtools}
\usepackage{algorithm}
\usepackage{algpseudocode}
\usepackage{geometry}
\geometry{margin=2.5cm}

\newcommand{\modtwelve}{\ (\mathrm{mod}\ 12)}
\newcommand{\omegaT}{\omega_{12}}

\title{\textbf{Messprotokoll: Phasenprogrammierung (mod 12) und Rekonstruktion von $\gamma$}}
\author{}
\date{}

\begin{document}
	\maketitle
	
	\section{Inputs und Konstanten}
	
	\paragraph{Inputs.}
	\begin{itemize}
		\item \textbf{Phasen-Lookup-Tabellen pro Kanal:} $\mathrm{DAC}_i(K)$ für $i\in\{E,A,B,C\}$ und $K\in\{0,\dots,11\}$.
		\item \textbf{Gewichte (Kalibration):} $w_E,w_A,w_B,w_C$ (Startwert z.\,B. $w_i\approx 1/4$).
		\item \textbf{Globaler Phasenoffset:} $\phi_0$ (aus Referenzmessung).
		\item \textbf{Messkampagnenliste:} Einträge $(K_E,K_A,K_B,R)$.
	\end{itemize}
	
	\paragraph{Konstante.}
	\[
	\omegaT \coloneqq e^{2\pi i/12}.
	\]
	
	\section{Definitionen: Mess- und Modellgröße}
	
	\paragraph{Gemessene komplexe Größe.}
	Aus den nichtselektiven Messungen von $\expval{\sigma_x}$ und $\expval{\sigma_y}$ auf dem System $S$ wird
	\[
	\gamma_{\mathrm{meas}}
	\;\coloneqq\;
	\bigl(s_x + i\,s_y\bigr)\,e^{-i\phi_0}
	\quad\text{mit}\quad
	s_x=\expval{\sigma_x},\;\; s_y=\expval{\sigma_y}.
	\]
	
	\paragraph{Vorhersage aus programmierten Phasen (Gewichtssuperposition).}
	Für die programmierte Wahl $(K_E,K_A,K_B,K_C)$ definiere
	\[
	\gamma_{\mathrm{pred}}
	\;\coloneqq\;
	w_E\,\omegaT^{K_E}
	+
	w_A\,\omegaT^{K_A}
	+
	w_B\,\omegaT^{K_B}
	+
	w_C\,\omegaT^{K_C}.
	\]
	
	\section{Kampagnen-Ablauf (Algorithmus)}
	
	\begin{algorithm}[h]
		\caption{Phasenprogrammierung (mod 12), Experiment, Rekonstruktion von $\gamma$}
		\begin{algorithmic}[1]
			\Require Lookup-Tabellen $\mathrm{DAC}_i(K)$; Gewichte $w_E,w_A,w_B,w_C$; Offset $\phi_0$; Kampagnenliste $(K_E,K_A,K_B,R)$
			\Ensure Datensätze $\{K_E,K_A,K_B,K_C,R,\gamma_{\mathrm{meas}},\gamma_{\mathrm{pred}}\}$
			
			\ForAll{Kampagneneinträge $(K_E,K_A,K_B,R)$}
			\State \textbf{(1) Bestimme $K_C$ durch Nebenbedingung}
			\State $K_C \gets (R - K_E - K_A - K_B)\modtwelve$
			\Statex
			
			\State \textbf{(2) Programmiere Phasen (über DAC)}
			\State $\texttt{set\_phase}(E,\,\mathrm{DAC}_E(K_E))$
			\State $\texttt{set\_phase}(A,\,\mathrm{DAC}_A(K_A))$
			\State $\texttt{set\_phase}(B,\,\mathrm{DAC}_B(K_B))$
			\State $\texttt{set\_phase}(C,\,\mathrm{DAC}_C(K_C))$
			\Statex
			
			\State \textbf{(3) Experiment ausführen}
			\State Präpariere Systemzustand $\ket{+}$
			\State Koppelung $S \leftrightarrow \text{Record}$ applizieren
			\State Record-Ausgänge ignorieren (nichtselektiv)
			\State Messe $s_x=\expval{\sigma_x}$ und $s_y=\expval{\sigma_y}$ auf $S$
			\State Speichere $(s_x,s_y)$
			\Statex
			
			\State \textbf{(4) Rekonstruiere $\gamma_{\mathrm{meas}}$ (global-phasen-korrigiert)}
			\State $\gamma_{\mathrm{meas}} \gets (s_x + i\,s_y)\,e^{-i\phi_0}$
			\Statex
			
			\State \textbf{(5) Berechne $\gamma_{\mathrm{pred}}$ aus $(K_E,K_A,K_B,K_C)$}
			\State $\gamma_{\mathrm{pred}} \gets w_E\omegaT^{K_E} + w_A\omegaT^{K_A} + w_B\omegaT^{K_B} + w_C\omegaT^{K_C}$
			\Statex
			
			\State \textbf{(6) Datensatz speichern}
			\State $\texttt{save}\{K_E,K_A,K_B,K_C,R,\gamma_{\mathrm{meas}},\gamma_{\mathrm{pred}}\}$
			\EndFor
		\end{algorithmic}
	\end{algorithm}
	
	\section{Postprocessing}
	
	\paragraph{(A) Komplexebene.}
	Plotte die Punkte $\gamma_{\mathrm{meas}}$ in der Ebene $(\Re,\Im)$, gruppiert nach
	\begin{itemize}
		\item festem $R$ (``$R$-Slices''), und/oder
		\item festen Tripeln $(K_E,K_A,K_B)$ (Trajektorien in $R$).
	\end{itemize}
	
	\paragraph{(B) Fehlermaß.}
	Definiere den Betragfehler
	\[
	\Delta \;\coloneqq\; \abs{\gamma_{\mathrm{meas}} - \gamma_{\mathrm{pred}}}.
	\]
	Optional zusätzlich: komponentenweise Fehler $(\Re,\Im)$ oder Winkel-/Radiusfehler.
	
	\paragraph{(C) Clustering.}
	Cluster $\gamma_{\mathrm{meas}}$ mittels Distanzschwelle $\varepsilon$ (aus Phasenjitter geschätzt),
	z.\,B. in der komplexen Ebene mit euklidischer Distanz.
	Berichte pro Cluster $c$:
	\[
	\mu_c = \text{Clusterzentrum},\qquad
	\Sigma_c = \text{Varianz/Kovarianz} \quad (\text{oder Radiusvarianz}).
	\]
	
	\paragraph{(D) Report.}
	\begin{itemize}
		\item Anzahl Cluster pro $R$.
		\item Clusterzentren und Varianzen pro $R$.
		\item $R$-Sweep-Trajektorien $\gamma(R)$ für feste Tripel $(K_E,K_A,K_B)$.
	\end{itemize}
	
	\paragraph{Nulltest (Minimale Amplitude und Muster).}
	Finde Einstellungen mit minimaler Amplitude
	\[
	\abs{\gamma_{\mathrm{meas}}}\ \text{minimal} \quad \text{(oder }\abs{\gamma_{\mathrm{pred}}}\text{ minimal).}
	\]
	Prüfe, ob diese Minimierer bevorzugt dem Muster
	\[
	\{u,\ v,\ u+6,\ v+6\}\modtwelve
	\]
	entsprechen (``Gegenphasen-Paarung'' um $6$).
	
\end{document}
